% =====================================================================
%  Fiche d’activités : les 4 activités du chapitre 5 (physique)
%  L’énergie thermique : le transfert thermique — deux pages A4
%  (activités 1 et 2 page 1, activités 3 et 4 page 2)
%  (sans situation de départ, sans zone de nom, sans bandeau de titre)
%  À coller dans le cahier. Compilation : xelatex fiche_unique.tex (deux passes)
% =====================================================================
\documentclass[10pt,a4paper]{article}
\usepackage[a4paper,top=.75cm,bottom=.6cm,left=.8cm,right=.8cm]{geometry}
\usepackage{fontspec}
\usepackage{polyglossia}
\setmainlanguage{french}
\usepackage{amsmath}
\usepackage{unicode-math}
\setmainfont{TeX Gyre Pagella}
\setsansfont{Poppins}[Path=fonts/, Extension=.ttf, UprightFont=*-Regular, BoldFont=*-Bold, ItalicFont=*-Italic, BoldItalicFont=*-BoldItalic]
\setmathfont{Latin Modern Math}
\setmathfont[range={it/{latin,Latin,greek,Greek}}]{texgyrepagella-italic.otf}
\setmathfont[range={up/{latin,Latin,num}}]{texgyrepagella-regular.otf}
\usepackage[table]{xcolor}
\usepackage{graphicx,tikz,array}
\usepackage[most]{tcolorbox}
\graphicspath{{images/}}
\usepackage{enumitem}
\usepackage{microtype}
\usepackage[newcommands]{ragged2e}
\setlength{\RaggedRightRightskip}{0pt plus 8em}
\usepackage[version=4]{mhchem}
\usepackage{siunitx}
\sisetup{locale=FR, per-mode=power, inter-unit-product=\ensuremath{\mathord{\cdot}}, group-minimum-digits=5}
\pagestyle{empty}
\setlength{\parindent}{0pt}
\setlength{\parskip}{0pt}
\relpenalty=10000 \binoppenalty=10000
\emergencystretch=3em

\definecolor{cInk}{HTML}{1B2A41}  \definecolor{cInk2}{HTML}{51627A}
\definecolor{cBlue}{HTML}{1F5FA8} \definecolor{cOrange}{HTML}{F07F1E}
\definecolor{cTeal}{HTML}{0E8F8A} \definecolor{cSoft}{HTML}{F4F7FB}
\definecolor{cRed}{HTML}{D62839}  \definecolor{cGreen}{HTML}{1F9D74}
\color{cInk}

\newcommand{\kw}[1]{{\sffamily\bfseries\color{cBlue}#1}}
\NewDocumentCommand{\docbadge}{m}{\tikz[baseline=(b.base)]\node[fill=cBlue,text=white,rounded corners=4pt,inner xsep=3pt,inner ysep=1.4pt,font=\sffamily\bfseries\tiny](b){Doc #1};}
% document : {largeur}{n}{légende}{image}
\newcommand{\doc}[4]{%
  \begin{minipage}[t]{#1}\vspace{0pt}\centering
    \includegraphics[width=\linewidth]{#4}\par\vspace{.5mm}
    {\sffamily\scriptsize\docbadge{#2}\enspace\color{cInk2}#3}
  \end{minipage}}
\NewDocumentCommand{\pastille}{m m}{\tikz[baseline=(n.base)]\node[circle,fill=#1,text=white,inner sep=0pt,minimum size=3.9mm,font=\sffamily\bfseries\tiny](n){#2};}
\setlist[enumerate,1]{label={\pastille{cBlue}{\arabic*}}, leftmargin=6mm, labelsep=1.8mm, itemsep=2pt, topsep=2pt, parsep=0pt}
% séparateur discret entre les activités
\newcommand{\sep}[1]{\par\vspace{2.5mm}%
  \noindent\tikz[baseline=(t.base)]\node[fill=cTeal,text=white,rounded corners=3pt,inner xsep=4pt,inner ysep=1.5pt,font=\sffamily\bfseries\scriptsize](t){Activité #1};%
  \hspace{2mm}{\color{cTeal!70}\leaders\hrule height .7pt\hfill\kern0pt}\par\vspace{1mm}}

% ---------- tableau de mesures compact ----------
\newcommand{\mesures}[3][0.8cm]{{\renewcommand{\arraystretch}{1.2}\setlength{\arrayrulewidth}{.5pt}\arrayrulecolor{cTeal!55}\setlength{\tabcolsep}{3pt}%
  \begin{tabular}{|>{\sffamily\scriptsize\bfseries\color{cTeal}}l|*{#2}{>{\centering\arraybackslash}p{#1}|}}\hline #3\end{tabular}}}

\newcommand{\vect}[1]{\overrightarrow{#1}}
\newcommand{\Ec}{E_{\mathrm{C}}}
\newcommand{\dC}[1]{\qty{#1}{\degreeCelsius}}
\newcommand{\cu}{\unit{J.kg^{-1}.\degreeCelsius^{-1}}}
\newcommand{\tfus}{\theta_{\mathrm{fus}}}
\begin{document}
\small\RaggedRight
\setlength{\parskip}{0pt}
% =====================================================================
\sep{1}
\begin{minipage}[t]{0.66\linewidth}\vspace{0pt}\RaggedRight
Des réchauds électriques identiques fournissent la même énergie par minute : l’énergie reçue est proportionnelle à la durée de chauffage. Les liquides sont pris à $\theta_i = \dC{20,0}$.\par\vspace{.5mm}
\textbf{Doc 1.} Pour élever la température de $\Delta\theta = \dC{15}$ : \qty{3,0}{min} pour \qty{100}{g} d’eau, \qty{6,0}{min} pour \qty{200}{g} d’eau.\par\vspace{.5mm}
\textbf{Doc 2.} \qty{200}{g} d’eau : durée de chauffage selon $\Delta\theta = \theta_f - \theta_i$.\enspace
\mesures[0.6cm]{3}{%
$\Delta\theta$ (\unit{\degreeCelsius}) & 10 & 20 & 30\\ \hline
Durée (\unit{min}) & 4,0 & 8,0 & 12,0\\ \hline}\par\vspace{.5mm}
\textbf{Doc 3.} Même durée de chauffage (\qty{4,0}{min}), donc même énergie, pour \qty{200}{g} d’eau et \qty{200}{g} d’huile.\enspace
\mesures[0.75cm]{2}{%
Liquide & eau & huile\\ \hline
$\Delta\theta$ (\unit{\degreeCelsius}) & 10,0 & 20,9\\ \hline}\par\vspace{.5mm}
\textbf{Information.} On note $Q$ la quantité de chaleur \emph{reçue} par un corps : $Q > 0$ s’il reçoit réellement de l’énergie, $Q < 0$ s’il en cède. Sans changement d’état, on admet que $Q$ est égale au produit de la masse $m$, de $\Delta\theta = \theta_f - \theta_i$ et d’un coefficient $c$ qui ne dépend que de la nature du corps : sa \textbf{chaleur massique}, en \cu{} (eau : 4\,185 ; huile : environ 2\,000). Le produit $C = m\cdot c$ est la \textbf{capacité thermique} du corps, en \unit{J.\degreeCelsius^{-1}}.\par\vspace{.5mm}
\begin{enumerate}
  \item \kw{Établir.} Docs 1 et 2 : montrer que la quantité de chaleur $Q$ reçue par l’eau est proportionnelle à sa masse $m$, puis à $\Delta\theta = \theta_f - \theta_i$.
  \item \kw{Interpréter.} Doc 3 : pour la même masse et la même énergie, l’élévation de température est-elle la même pour l’eau et l’huile ? Expliquer à l’aide des chaleurs massiques de l’information.
  \item \kw{Écrire.} À l’aide de l’information, écrire l’expression de $Q$ échangée par un corps pur de masse $m$ dont la température passe de $\theta_i$ à $\theta_f$, sans changement d’état. Préciser les unités.
  \item \kw{Définir.} Que représente $c$ quand $m = \qty{1}{kg}$ et $\Delta\theta = \dC{1}$ ? Que représente $C = m\cdot c$ ? Quel est le signe de $Q$ quand le corps se refroidit, et que signifie-t-il ?
\end{enumerate}
\end{minipage}\hfill
\begin{minipage}[t]{0.32\linewidth}\vspace{0pt}\centering
\doc{\linewidth}{1}{Deux masses d’eau}{doc1_chauffage}\par\vspace{1mm}
\doc{\linewidth}{3}{Eau et huile}{doc3_eau_huile}
\end{minipage}

% =====================================================================
\sep{2}
\begin{minipage}[t]{0.74\linewidth}\vspace{0pt}\RaggedRight
\textbf{Doc 4. Le calorimètre.} Enceinte thermiquement isolée : les corps qu’il contient échangent de la chaleur jusqu’à l’\textbf{équilibre thermique}, où ils ont tous la même température, constante. L’énergie cédée par les corps qui se refroidissent est entièrement reçue par ceux qui s’échauffent ; le bilan ne dépend que de l’état initial et de l’état final. Le calorimètre et ses accessoires reçoivent $Q = \mu\cdot(\theta_f - \theta_i)$, où $\mu$ est leur \textbf{capacité thermique}.\par\vspace{.4mm}
\textbf{Protocole 1.} $m_1 = \qty{200}{g}$ d’eau froide en équilibre à $\theta_1 = \dC{18,0}$ dans le calorimètre ; on ajoute $m_2 = \qty{200}{g}$ d’eau tiède à $\theta_2 = \dC{42,5}$ ; équilibre : $\theta_e = \dC{29,5}$.\par\vspace{.4mm}
\textbf{Protocole 2.} Un cylindre d’aluminium ($m_m = \qty{150}{g}$, $\theta_m = \dC{100,0}$) est introduit dans le même calorimètre contenant $m_e = \qty{250}{g}$ d’eau à $\theta_0 = \dC{20,0}$ ; équilibre : $\theta'_e = \dC{28,3}$. $c_e = \qty{4185}{J.kg^{-1}.\degreeCelsius^{-1}}$.\par\vspace{.5mm}
\begin{enumerate}
  \item \kw{Définir.} Quand des corps sont-ils en équilibre thermique ? D’après le doc 4, que vaut la somme $Q_1 + Q_2 + \dots$ des chaleurs échangées par tous les corps du calorimètre (\textbf{équation calorimétrique}) ? De quels états dépend le bilan ?
  \item \kw{Exprimer.} Protocole 1 : exprimer la chaleur $Q_1$ reçue par le calorimètre et l’eau froide, et la chaleur $Q_2$ cédée par l’eau tiède.
  \item \kw{Calculer.} Calculer la capacité thermique $\mu$ du calorimètre.
  \item \kw{Déterminer.} Protocole 2 : faire le bilan des chaleurs échangées, puis calculer la chaleur massique $c_m$ de l’aluminium (valeur admise : \qty{897}{J.kg^{-1}.\degreeCelsius^{-1}}).
\end{enumerate}
\end{minipage}\hfill
\doc{0.24\linewidth}{4}{Calorimètre (coupe)}{doc4_calorimetre}

\newpage
% =====================================================================
\sep{3}
\begin{minipage}[t]{0.66\linewidth}\vspace{0pt}\RaggedRight
\textbf{Doc 5.} On chauffe à puissance constante \qty{100}{g} de glace prise à \dC{-10} et on relève sa température.\par\vspace{.4mm}
\textbf{Chaleur latente.} À pression constante, un corps pur change d’état à température constante (\dC{0} pour la fusion de la glace, $\tfus$). La chaleur échangée est proportionnelle à la masse $m$ qui change d’état : $Q = m\cdot L$, $L$ étant la \textbf{chaleur latente}. La fusion se fait en recevant de la chaleur ; la solidification, inverse, en cède autant.\par\vspace{.4mm}
\textbf{Doc 6.} Calorimètre : $\mu = \qty{80}{J.\degreeCelsius^{-1}}$, $m_1 = \qty{300}{g}$ d’eau à $\theta_1 = \dC{25,0}$. On y introduit un glaçon de masse $m_2 = \qty{40}{g}$ à $\tfus = \dC{0,0}$. Après fusion complète : $\theta_e = \dC{13,3}$. $c_e = \qty{4185}{J.kg^{-1}.\degreeCelsius^{-1}}$.\par\vspace{.5mm}
\begin{enumerate}
  \item \kw{Exploiter.} Doc 5 : décrire les trois parties de la courbe. La glace reçoit-elle de l’énergie pendant le palier ?
  \item \kw{Définir.} D’après $Q = m\cdot L$, que représente $L$ quand $m = \qty{1}{kg}$ ? Quelle est son unité ? Quel est le signe de $Q$ pour une fusion, pour une solidification ?
  \item \kw{Décrire.} Doc 6 : décrire les deux étapes subies par le glaçon ; exprimer la chaleur totale $Q_g$ qu’il reçoit.
  \item \kw{Exprimer.} Exprimer la chaleur $Q_{ce}$ cédée par le calorimètre et l’eau.
  \item \kw{Calculer.} Calculer la chaleur latente de fusion de la glace $L_f$.
\end{enumerate}
\end{minipage}\hfill
\begin{minipage}[t]{0.32\linewidth}\vspace{0pt}\centering
\doc{\linewidth}{5}{Glace chauffée}{doc5_courbe}\par\vspace{1mm}
\doc{0.7\linewidth}{6}{Glaçon dans le calorimètre}{doc6_glacon}
\end{minipage}

% =====================================================================
\sep{4}
\begin{minipage}[t]{0.66\linewidth}\vspace{0pt}\RaggedRight
\textbf{Doc 7.} Un thermomètre au réservoir noirci est placé sous une cloche où l’on a fait le vide. Une lampe infrarouge, à l’extérieur, est dirigée vers le réservoir : la température augmente rapidement. La lampe émet lumière et infrarouge, des \textbf{ondes électromagnétiques} qui transportent de l’énergie, même dans le vide : c’est le \textbf{rayonnement}.\par\vspace{.4mm}
\textbf{Doc 8.} Une plaque noire mate et une plaque blanche brillante, identiques, sont à la même distance d’une lampe. Une surface \textbf{absorbe} une partie du rayonnement reçu (qui l’échauffe) et \textbf{réfléchit} le reste :\par\vspace{.4mm}
{\centering\mesures[0.6cm]{6}{%
$t$ (\unit{min}) & 0 & 2 & 4 & 6 & 8 & 10\\ \hline
$\theta_{\text{noire}}$ (\unit{\degreeCelsius}) & 20,0 & 28,5 & 34,6 & 39,0 & 42,1 & 44,3\\ \hline
$\theta_{\text{blanche}}$ (\unit{\degreeCelsius}) & 20,0 & 22,6 & 24,4 & 25,7 & 26,6 & 27,3\\ \hline}\par}\vspace{.5mm}
\textbf{Doc 9.} Dans un solide, la chaleur se propage de proche en proche, sans déplacement de matière : c’est la \textbf{conduction}. Dans un liquide ou un gaz, la matière chaude se déplace et transporte l’énergie : c’est la \textbf{convection}. Ces deux modes ont besoin de matière.\par\vspace{.5mm}
\begin{enumerate}
  \item \kw{Expliquer.} Doc 7 : pourquoi l’énergie ne peut-elle pas arriver au thermomètre par l’air ? Par quel mode de transfert le thermomètre reçoit-il de l’énergie ?
  \item \kw{Comparer.} Doc 8 : comparer l’échauffement des deux plaques. Laquelle absorbe le plus le rayonnement ? Laquelle le réfléchit le plus ?
  \item \kw{Distinguer.} Citer des exemples de rayonnement (chauffage de la Terre par le Soleil, chauffe-eau solaires…). D’après le doc 9, quelle différence essentielle y a-t-il entre le rayonnement et la conduction ou la convection ?
\end{enumerate}
\end{minipage}\hfill
\begin{minipage}[t]{0.32\linewidth}\vspace{0pt}\centering
\doc{\linewidth}{7}{Thermomètre sous vide}{doc7_cloche}\par\vspace{1mm}
\doc{\linewidth}{8}{Plaque noire et plaque blanche}{doc8_plaques}
\end{minipage}
\end{document}
